\section*{Formal author conditional algebraic-truncation results}
\paragraph{Theorem 32.} Let $k\in\mathbb N$. If $\tau(\mathrm{trunc}((1+k^2x)^{i/k},d))=O(\log^cd)$ (for some constant $c$) for all $i\in[k-1]$ then integer factorization (in the non-uniform setting) can be performed in polynomial time.
\paragraph{Theorem 35.} Fix $k\in\mathbb N$. Suppose, for each $i\in[k-1]$, there exists some constant $c$ such that $\tau(\mathrm{trunc}((1+k^2x)^{i/k},d))=O(\log^cd)$, for large enough $d$. Then, $(n!)_{n\in\mathbb N}$ is ultimately easy.
